% !TEX TS-program = lualatexmk
% !TEX parameter =  --shell-escape

\documentclass[vespers_book_la_eng.tex]{subfiles}

\ifcsname preamble@file\endcsname
  \setcounter{page}{\getpagerefnumber{M-vespers_book_la_eng_sundays_lent}}
\fi

\begin{document}

 \chapter*{SUNDAYS OF LENT AND OF PASSIONTIDE.}
\addcontentsline{toc}{chapter}{Sundays of Lent and of Passiontide.}\phantomsection

\header{sundays of lent at vespers.}

\bigtitle{I Sunday of Lent.}
\addcontentsline{toc}{section}{I Sunday of Lent.} \phantomsection

\chapterlateng

\scripture{2 Cor. 6, 1 – 2.}
\texteparacol{\lettrine{F}{r}atres: Hortámur vos, ne in vácuum grátiam Dei recipiátis.~† Ait enim: Témpore accépto exaudívi te,~* et in die salútis adjúvi te.}{english}{\noindent Brethren, we exhort you, that you receive not the grace of God in vain. For he saith: In an accepted time have I heard thee; and in the day of salvation have I helped thee.}

\hymnlateng

\label{hy_audi_benigne_conditor_solesmes_1961} \phantomsection

\gscore[]{2.}{hy_audi_benigne_conditor_solesmes_1961}

\textes{audi_benigne_en}

\label{versiculus_quadragesimae} \phantomsection

\smallscore{versiculus_quadragesimae}

\begin{enpars}\noindent \vv  God hath given his Angels charge over thee.

\noindent \rr To keep thee in all thy ways.\end{enpars}

\gscore[Ad Magnif.]{Ant. 8. G*}{an_ecce_nunc_tempus_solesmes_1961}
\antiphona{english}{Behold, now is the accepted time behold, now is the day of salvation; in these days therefore let us approve ourselves as the ministers of God, in much patience, in fastings, in watchings, and in love unfeigned}

\collectlateng

\texteparacol{\lettrine{D}{e}us, qui Ecclésiam tuam ánnua quadragesimáli observatióne puríficas:~† præsta famíliæ tuæ; ut quod a te obtinére abstinéndo nítitur,~* hoc bonis opéribus exsequátur. Per Dóminum.}{english}
{\noindent O God, who dost every year purge thy church by the fast of forty days, grant unto this thy family, that what things soever they strive to obtain at thy hand by abstaining from meats, they may ever turn to profit by good works.}

\bigtitle{II Sunday of Lent.}
\addcontentsline{toc}{section}{II Sunday of Lent.} \phantomsection

\chapterlateng

\texteparacol{\lettrine{F}{r}atres: Rogámus vos, et obsecrámus in Dómino Jesu:~† ut, quemádmodum accepístis a nobis, quómodo vos opórteat ambuláre, et placére Deo,~* sic et ambulétis, ut abundétis magis.}{english}{\noindent For the rest therefore, brethren, we pray and beseech you in the Lord Jesus, that as you have received from us, how you ought to walk, and to please God, so also you would walk, that you may abound the more.}

\rubrique{Hymnus \normaltext{Audi benígne Cónditor, \pageref{hy_audi_benigne_conditor_solesmes_1961}.} \normaltext{℣. Angelis suis, \pageref{versiculus_quadragesimae}.}}

\gscore[Ad Magnif.]{Ant. 1. f}{an_visionem_solesmes_1961}
\antiphona{english}{Tell the vision that ye have seen to no man, until the Son of man be risen again from the dead}

\collectlateng

\texteparacol{\lettrine{D}{e}us, qui cónspicis omni nos virtúte destítui: intérius exteriúsque custódi;~† ut ab ómnibus adversitátibus muniámur in córpore,~* et a právis cogitatiónibus mundémur in mente. Per Dóminum.}{english}
{\noindent O God, Who seest that we have no power of ourselves to help ourselves, keep us both outwardly in our bodies, and inwardly in our souls, that we may be defended from all adversities which may happen to the body, and from all evil thoughts which may assault and hurt the soul.}

\bigtitle{III Sunday of Lent.}
\addcontentsline{toc}{section}{III Sunday of Lent.} \phantomsection

\chapterlateng

\scripture{Eph. 5, 1 – 2.}
\texteparacol{\lettrine{F}{r}atres: Estóte imitatóres Dei, sicut fílii caríssimi:~† et ambuláte in dilectióne, sicut et Christus diléxit nos, et trádidit semetípsum pro nobis~* oblatiónem et hóstiam Deo in odórem suavitátis.}{english}{\noindent Brethren: Be ye, therefore, followers of God, as most dear children; And walk in love, as Christ also hath loved us, and hath delivered himself for us, an oblation and a sacrifice to God for an odour of sweetness.}

\rubrique{Hymnus \normaltext{Audi benígne Cónditor, \pageref{hy_audi_benigne_conditor_solesmes_1961}.} \normaltext{℣. Angelis suis, \pageref{versiculus_quadragesimae}.}}

\gscore[Ad Magnif.]{Ant. 8. G}{an_extollens_solesmes_1961}
\antiphona{english}{A certain woman of the company lifted up her voice and said: Blessed is the womb that bare thee, and the paps which Thou hast sucked. But Jesus said unto her: Yea, rather, blessed are they that hear the word of God, and keep it}

\collectlateng

\texteparacol{ \lettrine[depth=1]{Q}{u}ǽsumus omnípotens Deus, vota humílium réspice:~† atque ad defensiónem nostram,~* déxteram tuæ majestátis exténde.
Per Dóminum.}{english}
{\noindent We beeseech thee, almighty God, look upon the hearty desires of thy humble servants, and stretch forth the right hand of thy majesty to be our defense against all our enemies.}

\bigtitle{IV Sunday of Lent.}
\addcontentsline{toc}{section}{IV Sunday of Lent.} \phantomsection

\chapterlateng

\scripture{Gal. 4, 22 – 24.}
\texteparacol{\lettrine{F}{r}atres: Scriptum est quóniam Abraham duos fílios hábuit: unum de ancílla, et unum de líbera:~† sed qui de ancílla, secúndum carnem natus est: qui autem de líbera, per repromissiónem:~* quæ sunt per allegoríam dicta.}{english}{\noindent Brethren: For it is written that Abraham had two sons: the one by a bondwoman, and the other by a free woman. But he who was of the bondwoman, was born according to the flesh: but he of the free woman, was by promise, which things are said by an allegory.}

\rubrique{Hymnus \normaltext{Audi benígne Cónditor, \pageref{hy_audi_benigne_conditor_solesmes_1961}.} \normaltext{℣. Angelis suis, \pageref{versiculus_quadragesimae}.}}

\gscore[Ad Magnif.]{Ant. 1. g}{an_subiit_ergo_solesmes_1961}
\antiphona{english}{And Jesus went up into a mountain, and there He sat with his disciples}

\collectlateng

\texteparacol{\lettrine{C}{o}ncéde, quǽsumus omnípotens Deus:~† ut, qui ex mérito nostræ actiónis afflígimur,~* tuæ grátiæ consolatióne respirémus.
Per Dóminum.}{english}
{\noindent Grant, we beseech thee, Almighty God, that we who for our evil deeds are worthily punished, may, by the comfort of thy grace, mercifully be relieved.}

\bigtitle{Passion Sunday.}
\addcontentsline{toc}{section}{Passion Sunday.} \phantomsection

\chapterlateng

\scripture{Heb. 9, 11 – 12.}

\texteparacol{\lettrine{F}{r}atres: Christus assístens Póntifex futurórum bonórum, per ámplius et perféctius tabernáculum non manu fa\-ctum, id est, non huius creatiónis:~† neque per sánguinem hircórum aut vitulórum, sed per próprium sánguinem introívit semel in San\-cta,~* ætérna redemptióne invénta.}{english}{\noindent Brethren: But Christ, being come an high priest of the good things to come, by a greater and more perfect tabernacle not made with hand, that is, not of this creation: Neither by the blood of goats, or of calves, but by his own blood, entered once into the holies, having obtained eternal redemption.}

\hymnlateng

\label{hy_vexilla_regis_prodeunt_solesmes} \phantomsection

\rubrique{\begin{enpars}All kneel for the entire strophe beginning \normaltext{O crux ave Spes única.}\end{enpars}}
 
 \gscore[]{1.}{hy_vexilla_regis_prodeunt_solesmes}
 
 \textes{vexilla_regis_en}
 
 \label{versiculus_passionis} \phantomsection
 \smallscore{versiculus_passionis}
 
 \begin{enpars}\noindent \vv Deliver me, O Lord, from the evil man.

\noindent \rr Rescue me from the unjust man.\end{enpars}

\gscore[Ad Magnif.]{Ant. 2. D}{an_abraham_pater_vester_solesmes_1961} 
\antiphona{english}{Your father Abraham rejoiced to see my day and he saw it, and was glad}

\collectlateng

\texteparacol{\lettrine[depth=1]{Q}{u}ǽsumus omnípotens Deus, famíliam tuam propítius réspice:~† ut te largiénte, regátur in córpore;~* et te servánte, custodiátur in mente. Per Dóminum.}{english}
{\noindent We beseech thee, Almighty God, mercifully to look upon this thy family, that by thy great goodness they may be governed and preserved evermore, both in body and soul. Through our Lord.}

\bigtitle{Palm Sunday.}
\addcontentsline{toc}{section}{Palm Sunday.} \phantomsection
\header{sundays of passiontide at vespers.}

\chapterlateng

\scripture{Phil. 2, 5 – 7.}
\texteparacol{\lettrine{F}{r}atres: Hoc enim sentíte in vobis, quod et in Christo Jesu: qui, cum in forma Dei esset, non rapínam arbitrátus est esse se æquálem Deo:~† sed semetípsum exinanívit, formam servi accípiens, in similitúdinem hóminum factus,~* et hábitu invéntus ut homo.}{english}{\noindent Brethren: For let this mind be in you, which was also in Christ Jesus: Who being in the form of God, thought it not robbery to be equal with God: But emptied himself, taking the form of a servant, being made in the likeness of men, and in habit found as a man.}

\rubrique{Hymnus \normaltext{Vexílla Regis pródeunt, \pageref{hy_vexilla_regis_prodeunt_solesmes}.} \normaltext{℣. Éripe, \pageref{versiculus_passionis}.}}

\gscore[Ad Magnif.]{Ant. 8. G*}{an_scriptum_est_enim__percutiam_solesmes_1961}

\antiphona{english}{It is written: I will smite the Shepherd, and the sheep of the flock shall be scattered abroad, but after I am risen again, I will go before you into Galilee: there shall ye see Me, saith the Lord}

\collectlateng

\texteparacol{\lettrine{O}{m}nípotens sempitérne Deus, qui humáno géneri, ad imitándum humilitátis exémplum, Salvatórem nostrum carnem súmere, et crucem subíre fecísti:~† concéde propítius; ut et patiéntiæ ipsíus habére documénta,~*  et resurrectiónis consórtia mereámur. Per eúmdem Dóminum nostrum.}{english}
{\noindent Almighty and everlasting God, who, of thy tender love towards mankind, hast sent thy Son our saviour Jesus Christ to take upon him our flesh and to suffer death upon the Cross, that all mankind should follow the example of his great humility; mercifully grant, that we may both follow the example of his patience, and also be made partakers of his resurrection. Through the same Lord.}

\end{document}